% 7.模型评价.tex —— 七、模型的评价
\section{模型的评价}

\subsection{模型的优点}
\begin{itemize}
  \item \textbf{优点 1}：数值方法递进清晰，先以 Euler 显式方法完整呈现中心点、内部点、边界点三类节点的离散构造，再引入 Crank--Nicolson 与 Thomas 追赶法，使格式的物理含义与升级动机一目了然。
  \item \textbf{优点 2}：中心点用 L'H\^opital 法则消去 $1/r$ 奇异性、边界点用半控制体积精确平衡，二者均保持二阶精度与离散守恒（守恒偏差 $10^{-14}$ 量级）。
  \item \textbf{优点 3}：格式对比量化，显式方法的 CFL 边界（$Fo_{\text{crit}}=0.413$）由特征值精确导出并与数值实验一致，CN 的无条件稳定性收益得到定量验证。
\end{itemize}

\subsection{模型的缺点}
\begin{itemize}
  \item \textbf{缺点 1}：未计蒸发潜热（题设未给出汽化潜热），热方程无汇项，会高估升温速率与温度水平，是本模型的主要系统误差来源。
  \item \textbf{缺点 2}：水分侧 Robin 边界条件为传质类比假设，$D(C)$ 经验公式的指数形式解读存在不确定性，二者直接影响浓度结果的量级。
\end{itemize}
